% SGA 6, Exposición XIII: fragmento español alineado con la fuente.
% Alcance exacto: fr.tex, líneas 14751--14950 inclusive.
% El límite final queda dentro del diagrama del Lema 3.4 y continúa en el
% fragmento siguiente; no se añaden cierres ajenos a la fuente.
% SHA-256 de la fuente francesa completa:
% 0D2158B9FE3B466309A92165727A0B191422E05FBAA0DACD5806E8322C54BA1A.

\textbf{Teorema 1.13.} \emph{Sean \(S\) un esquema noetheriano y \(X\) un
\(S\)-esquema proyectivo, provisto de un haz amplio \(\mathcal O_X(1)\)
cuyos inducidos \(\mathcal O_{X_s}(1)\), con \(s\in S\), están generados por
su \(H^0(\mathcal O_{X_s}(1))\); sea \(\mathcal F\) una familia de clases de
haces coherentes sobre las fibras de \(X/S\). Entonces las condiciones
siguientes son equivalentes}:
\begin{enumerate}[label=(\roman*)]
\item \emph{\(\mathcal F\) es acotada. Además, si todos los
\(F_K\in\mathcal F\) son localmente libres de rango \(p\), se puede suponer
que \(\mathcal F\) está acotada por un haz \(F\) sobre un \(X_T\), con \(F\)
localmente libre de rango \(p\).}
\item \emph{El conjunto de los polinomios de Hilbert \(\chi(F_K(n))\) de los
\(F_K\in\mathcal F\) es finito, y existe una sucesión de enteros \((b)\) tal
que \(\mathcal F\) sea una \((b)\)-familia.}
\item \emph{El conjunto de los polinomios de Hilbert \(\chi(F_K(n))\) de los
\(F_K\in\mathcal F\) es finito, y existe un entero \(m\) tal que todos los
\(F_K\in\mathcal F\) sean \(m\)-regulares.}
\item \emph{El conjunto de los polinomios de Hilbert \(\chi(F_K(n))\) de los
\(F_K\in\mathcal F\) es finito, y \(\mathcal F\) está contenida en la
familia de las clases de cocientes de haces de la forma \(E_K\), donde \(E\)
es un haz coherente sobre un \(X_T\). Además, se puede suponer \(T=S\) y
\(E\) de la forma \(\mathcal O_X(-m)^{\oplus M}\).}
\item \emph{\(\mathcal F\) está contenida en la familia de las clases de
conúcleos de homomorfismos de la forma \(E'_K\to E_K\), donde \(E'\) y \(E\)
son haces coherentes sobre un \(X_T\). Además, se puede suponer \(T=S\), y
\(E'\) y \(E\) de las formas \(\mathcal O_X(-m')^{\oplus M'}\) y
\(\mathcal O_X(-m)^{\oplus M}\).}
\end{enumerate}
En efecto, supongamos que \(\mathcal F\) está acotada por un haz \(F\) sobre
un \(X_T\); descomponiendo \(T\) en partes si es necesario y utilizando el
teorema de planitud genérica (EGA IV 6.9.1), se puede suponer que \(F\) es
plano sobre \(T\). Entonces, el número de polinomios \(\chi(\sisi{F}{F_t}(n))\) es a lo % Canon SGA-SRC-001425
sumo igual al número de componentes conexas de \(T\) (EGA III 7.9.5), y un
pequeño lema afirma que, si para un \(t\in T\), \(F_t\) es localmente libre
de rango \(p\), entonces \(F\) también lo es por encima de un entorno
abierto de \(t\) en \(T\) (véase EGA IV 11.3.10, donde se toma \(X=Y\),
\(f=\mathrm{id}\)). Ahora bien, de las hipótesis y de EGA III 7.7.6 resulta
que, descomponiendo \(T\) todavía más finamente si es necesario, se puede
encontrar una sucesión \(\sigma\) de secciones
\(\sigma_1,\ldots,\sigma_r\in H^0(\mathcal O_X(1))\) que es
\(F\)-regular. Por tanto, (ii) resulta de la propiedad de semicontinuidad
superior de la función \(t\longmapsto h^0(F_{t,i}(-1))\), donde, para
\(t\in T\) y \(0\le i\le r\), \(F_{t,i}\) designa la restricción de \(F_t\)
al esquema de ceros de \(\sigma_1,\ldots,\sigma_i\) (EGA III 7.7.5).

La implicación (ii) \(\Rightarrow\) (iii) resulta inmediatamente de 1.11.
La implicación (iii) \(\Rightarrow\) (iv) resulta de 1.3 (iii), tomando
\(M=\max(\chi(F_K(m)))\) y \(E=\mathcal O_X(-m)^{\oplus M}\).

Supongamos que \(\mathcal F\) satisface (iv); existen sucesiones exactas
\[
0\longrightarrow F'_K\longrightarrow E_K\longrightarrow F_K\longrightarrow0
\]
y el conjunto de los polinomios \(\chi(F_K(n))\) es finito. Por hipótesis,
la familia de las clases de los \(E_K\) es acotada; por la implicación (i)
\(\Rightarrow\) (ii), ya demostrada, el conjunto de los polinomios
\(\chi(E_K(n))\) es, pues, finito, y existe una sucesión de enteros \((b)\)
tal que todos los \(E_K\) (con \(K\) algebraicamente cerrado) sean
\((b)\)-haces. Por consiguiente, el conjunto de los polinomios
\(\chi(F'_K(n))\) es finito y, en virtud de 1.6, todos los \(F'_K\) son
\((b)\)-haces. Aplicando entonces a la familia de las clases de los \(F'_K\)
la implicación (ii) \(\Rightarrow\) (iv), ya demostrada, se obtiene (v).

Supongamos finalmente que \(\mathcal F\) satisface (v), y demostremos que
\(\mathcal F\) es acotada. Se puede suponer que \(E\) (resp. \(E'\)) es de
la forma \(\mathcal O_X(-m)^{\oplus M}\). En efecto, por (i)
\(\Rightarrow\) (iv), tomemos una sobreyección
\(L=\mathcal O_{X_T}(-m)^{\oplus M}\longrightarrow E\). Descomponiendo \(T\)
en partes si es necesario, de modo que \(E\) pase a ser plano sobre \(T\),
la formación de la sucesión exacta
\(0\longrightarrow I\xrightarrow{u}L\longrightarrow E\longrightarrow0\)
conmuta con el paso a las fibras. Por (i) \(\Rightarrow\) (iii) y (i)
\(\Rightarrow\) (iv), tomemos \(m_1\) suficientemente grande para que todos
los
\(\operatorname{Ext}^1(\mathcal O_{X_K}(-m_1),I_K)=H^1(I_K(m_1))\)
sean nulos y para que existan sobreyecciones
\(L_1=\mathcal O_{X_T}(-m_1)^{\oplus M_1}\longrightarrow E'\) y
\(\alpha:L_2=\mathcal O_{X_T}(-m_1)^{\oplus M_2}\longrightarrow I\).
Entonces las aplicaciones
\(\operatorname{Hom}(L_{1,K},L_K)\longrightarrow
\operatorname{Hom}(L_{1,K},E_K)\) son sobreyectivas. Sea
\(\beta:E'_K\to E_K\) un homomorfismo y sea \(\gamma\) la composición
\(\sisi{L'_K}{L_{1,K}}\to E'_K\to E_K\); entonces \(\gamma\) proviene de un homomorfismo % Canon SGA-SRC-001426
\(\delta:L_{1,K}\to L_K\), y
$(\delta,u_K\circ\alpha):L_{1,K}\oplus L_{2,K}\longrightarrow L_K$ y
\(\beta:E'_K\to E_K\) tienen el mismo conúcleo. Por consiguiente,
\(\mathcal F\) está contenida en la familia de las clases de los conúcleos de
homomorfismos de la forma
\(\mathcal O_{X_K}(-m_1)^{\oplus(M_1+M_2)}\longrightarrow
\mathcal O_{X_K}(-m)^{\oplus M}\).

Descomponiendo siempre \(T\) en partes si es necesario, se puede suponer
además que \(X_T\) (y, por tanto,
\(\mathcal H\!om_{\mathcal O_{X_T}}(E',E)\)) y todos los
\(R^q(f_T)_*\mathcal H\!om_{\mathcal O_{X_T}}(E',E)\) (donde
\(f_T:X_T\to T\) es el morfismo estructural) son planos sobre \(T\), a fin de
que la formación de
\(G=(f_T)_*\mathcal H\!om_{\mathcal O_{X_T}}(E',E)\) conmute con cualquier
cambio de base \(T'\to T\) (EGA III 6.9.9.2). Sea $R=\mathbb V(G^*)$; entonces
\(X_R\) lleva canónicamente una sucesión exacta «~universal~»
\[
E'_R\longrightarrow E_R\longrightarrow F\longrightarrow0
\]
y \(F\) acota \(\mathcal F\).

\subsection*{2. Varios lemas técnicos}
\textbf{2.1.} Sean \(k\) un cuerpo y \(X\) un \(k\)-esquema propio. Sean
\(L_1,\ldots,L_s\) haces invertibles sobre \(X\), \(c_1(L_i)\) la clase de
Chern de \(L_i\), e \(Y\) un subesquema cerrado de \(X\) de dimensión
\(\le s\). Recordemos que se designa por
\(\langle c_1(L_1)\cdots c_1(L_s)\cdot Y\rangle\) el número de intersección,
que puede calcularse como el coeficiente del monomio \(n_1\cdots n_s\) en el
polinomio de Snapper
\(\chi\bigl(L_1^{\otimes n_1}\otimes\cdots\otimes L_s^{\otimes n_s}\otimes\mathcal O_Y\bigr)\)
(X 4.3.1).

\textbf{2.2.} En adelante, se supone que \(k\) es algebraicamente cerrado,
que \(X\) es proyectivo, provisto de un haz invertible muy amplio
\(H=\mathcal O_X(1)\), y que \(X\) es íntegro, de dimensión \(r\).

Sea \(L\) un haz invertible sobre \(X\). El entero
\(\langle c_1(L)\cdot c_1(H)^{r-1}\rangle\) se denomina el \emph{grado} de
\(L\) (relativo a \(H\)) y se denota por \(d(L)\).

En lo que sigue, se considera un haz invertible \(L\) y los polinomios de
Hilbert: \(\chi(L(n))=\sum_{i=0}^{r}a_i\binom{n+i}{i}\) y
\(\chi(\mathcal O_X(n))=\sum_{i=0}^{r}e_i\binom{n+i}{i}\).

\textbf{Lema 2.3.} \emph{Se tiene}:
\begin{enumerate}[label=(\roman*)]
\item \emph{\(a_r=e_r=d(H)\).}
\item \emph{(Si \(r\ge1\)), \(a_{r-1}-e_{r-1}=d(L)\).}
\end{enumerate}

En efecto, se obtienen (i) y (ii) inmediatamente tomando \(m=0,1\) en el
polinomio siguiente:
\[
\chi\bigl(L^{\otimes m}\otimes H^{\otimes n}\bigr)
=
\sum_{i+j\le r} f_{ij}\binom{m+i}{i}\binom{n+j}{j}.
\]

\textbf{Lema 2.4.} \emph{Si \(d(L)\le0\) y \(h^0(L)\ne0\), entonces
\(L\simeq\mathcal O_X\); por tanto, \(d(L)=0\) y \(h^0(L)=1\).}

En efecto, como \(X\) es íntegro, existe un divisor \(Y\ge0\) tal que
\(L\simeq\mathcal O_X(Y)\). Si \(Y>0\), se tendría
\(0\ge\langle \sisi{c(L)}{c_1(L)}\cdot c_1(H)^{r-1}\rangle=\langle c_1(H)^{r-1}\cdot Y\rangle\ge1\). % Canon SGA-SRC-001427

\textbf{Lema 2.5.} \emph{Si \(d(L)\le0\), entonces \(L\) es un
\((0,\ldots,0,a_r)\)-haz (con \(r-1\) ceros).}

En efecto, si \(r=0\), \(L\) es evidentemente un \(a_0\)-haz. Si \(r\ge1\),
por 2.4, \(h^0(L(-1))=0\), puesto que \(d(L(-1))<0\). Sea \(Y\) una sección
hiperplana de \(X\), y sea
\(0\longrightarrow L(-1)\longrightarrow L\longrightarrow L_Y\longrightarrow0\)
una sucesión exacta definida por \(Y\). Si \(r=1\), por 1.7,
\(h^0(L_Y(-1))=a_1\); por tanto, \(L\) es un \((0,a_1)\)-haz. Si \(r\ge2\),
se puede tomar \(Y\) íntegro gracias al teorema de Bertini, y se puede
suponer así, por inducción, que \(L_Y\) es un
\((0,\ldots,a_r)\)-haz; de donde resulta la afirmación.

\textbf{Lema \sisi{2.7}{2.6}.} \emph{Si \(r=2\) y \(d(L)\le0\), entonces}: % Canon SGA-SRC-001428
\begin{enumerate}[label=(\roman*)]
\item \emph{Para todo \(n\), \(h^1(L(n))\le a_1(a_1+1)a_2-a_0\)}
\item \emph{\(a_0\le a_1(a_1+1)\).}
\end{enumerate}

Sea \(Y\) una sección hiperplana íntegra de \(X\). Por 1.7 y 2.6, \(L_Y\)
es un \((0,a_2)\)-haz con polinomio de Hilbert
\(\chi(L_Y(n))=a_1+a_2\binom{n+1}{1}\). Por tanto, por 1.11,
\(-a_1\ge0\) y \(L_Y\) es \([-a_1]\)-regular.

Consideremos la sucesión exacta siguiente:
\[
H^0(L_Y(n))\longrightarrow H^1(L(n-1))
\longrightarrow H^1(L(n))\longrightarrow H^1(L_Y(n)).
\]
Para \(n\ge -a_1-1\), se tiene
\(h^1(L(-a_1-2))\ge h^1(L(n))\), pues \(h^1(L_Y(n))=0\) por 1.3 (i); para
\(n\le -1\), se tiene \(h^1(L(n-1))\le h^1(L(-1))\), pues
\(h^0(L_Y(n))=0\) por 2.4. Finalmente, para todo \(n\),
\(h^0(L_Y(n-1))\le h^0(L_Y(n))\), pues existe una inyección
\(L_Y(n-1)\longrightarrow L_Y(n)\), y
\(h^0(L_Y(-a_1-1))=a_2\binom{-a_1}{1}+a_1\), pues
\(h^1(L_Y(-a_1-1))=0\); por tanto, para \(n\le -a_1-2\), se tiene
\(h^1(L(n))-h^1(L(n+1))\le -a_2a_1+a_1\). Como consecuencia, se concluye
que, para todo \(n\),
\[
h^1(L(n))\le h^1(L(-a_1-2))+(-a_1-1)(-a_2a_1+a_1).
\tag{\sisi{2.7.1}{2.6.1}}
\]

Como \(L_Y\) es \([-a_1]\)-regular, \(h^2(L(-a_1-2)\sisi{}{)}=0\) por 1.4; de donde % Canon SGA-SRC-001429
\(h^1(L(-a_1-2))=h^0(L(-a_1-2))-\chi(L(-a_1-2))\). Ahora bien,
\(h^0(L(-a_1-2))\le a_2\binom{-a_1}{2}\) por 1.8 (i), ya que \(L\) es un
\((0,0,a_2)\)-haz por 2.5. Por consiguiente,
\[
0\le h^1(L(-a_1-2))\le -a_1(-a_1-1)-a_0.
\tag{2.7.2}
\]
Finalmente, 2.7.2 da (ii), y 2.7.1 y 2.7.2 dan (i).

\textbf{Lema 2.8.} \emph{Supongamos que \(r=2\) y \(d(L)\le0\). Sean
\(\delta=d(L)\) y \(\beta=\delta(e_2-1)+e_1\). Entonces}:
\[
a_0\ge
\langle c_1(L)^2\rangle+2e_0+
\sum_{i=1}^{2}(2e_i-a_i)\binom{\delta-1+i}{i}
+\beta(\beta+1).
\]

En efecto, sea
\(\chi\bigl(L^{\otimes m}(\delta-1)\bigr)=\alpha_2m^2+\alpha_1m+\alpha_0\);
entonces, comparando polinomios, resulta que
\(2\alpha_2=\langle c_1(L)^2\rangle\), que
\(\alpha_0=\sum_{i=0}^{2}e_i\binom{\delta-1+i}{i}\) y que
\(\alpha_2+\alpha_1+\alpha_0=\sum_{i=0}^{2}a_i\binom{\delta-1+i}{i}\).
Sea
\(\chi\bigl(L^{-1}(\delta)(n)\bigr)=\sum_{i=0}^{2}\beta_i\binom{n+i}{i}\);
entonces \(\beta_0=\alpha_2-\alpha_1+\alpha_0\) y, por 2.3,
\(\beta_1=d(L^{-1}(\delta))+e_1\), lo que es igual a \(\beta\). Por tanto,
\(a_0=2\alpha_2+2\alpha_0-\sum_{i=1}^{2}a_i\binom{\delta-1+i}{i}-\beta_0\),
y, como \(d(L^{-1}(\delta))\le0\), se tiene
\(\beta_0\le\beta(\beta+1)\) por 2.7 (ii); de donde resulta la afirmación.

\textbf{Lema 2.9.} \emph{Supongamos que \(r\ge2\), que \(d(L)\le0\) y que
\(X\) es normal. Entonces}:
\begin{enumerate}[label=(\roman*)]
\item \emph{\(h^1(L(-n))=0\) para
\(n\ge2+h_2(a_{r-2},a_{r-1},a_r)\),}
\item \emph{\(h^1(L(n))\le h_r(a_{r-2},a_{r-1},a_r)\) para todo \(n\),}
\item \emph{\(|a_i|\le g_{r-i}(a_{r-2},a_{r-1},a_r)\) para
\(i=0,\ldots,r\),}
\end{enumerate}
\emph{donde
\(h_2(X_0,X_1,X_2),h_3(X_0,X_1,X_2),\ldots\) y
\(g_0(X_0,X_1,X_2),g_1(X_0,X_1,X_2),\ldots\) son polinomios universales con
coeficientes en \(\mathbb Q\), que no dependen ni de \(k\) ni de \(X\).}

En efecto, tratemos primero el caso \(r=2\). Sean
\(g_i(X_0,X_1,X_2)=X_{2-i}^2\) para \(i=0,1,2\) y
\(h_2(X_0,X_1,X_2)=X_1(X_1+1)X_2-X_0\). Entonces (iii) es evidente y (ii)
es 2.7 (i).

Sean \(X\) en \(P=\mathbb P_k^N\) y
\(\omega=\mathcal E\!xt^{N-2}_{\mathcal O_P}(\mathcal O_X,\Omega^N_{P/k})\).
Entonces \(h^1(L(-n))=h^1(\omega\otimes L^{-1}(n))\), y \(\omega\) es
invertible fuera del conjunto finito \(Z\) de los puntos singulares de \(X\)
(véase FGA 149 \S 6). Sea \(Y\) una sección hiperplana normal de \(X\) que
no contiene a \(Z\). Entonces \(\omega_Y(1)\simeq\Omega^1_{Y/k}\), y se tiene
una sucesión exacta:
\[
0\longrightarrow \omega\otimes L^{-1}(-1)
\longrightarrow \omega\otimes L^{-1}
\longrightarrow \Omega^1_{Y/k}\otimes L_Y^{-1}(-1)
\longrightarrow0.
\]
Ahora bien, \(h^1(\Omega^1_{Y/k}\otimes L_Y^{-1}(n))=h^0(L_Y(-n))\) para
todo \(n\), y \(h^0(L_Y(-n))=0\) para \(n\ge1\) por 2.5. Por tanto,
\(\Omega^1_{Y/k}\otimes L_Y^{-1}(-1)\) es 3-regular y, por 1.4,
\(\omega\otimes L^{-1}\) es
\(\bigl[3+h^1(\omega\otimes L^{-1}(2))\bigr]\)-regular. En consecuencia,
\(h^1(L(-n))=0\) para \(n\ge2+h^1(L(-2))\)\,; de donde resulta (i).

Supongamos ahora que \(r\ge3\) y razonemos por inducción sobre \(r\). Sea
\(Y\) una sección hiperplana normal de \(X\), y consideremos la sucesión
exacta:
\[
H^0(L_Y(-n))\longrightarrow H^1(L(-n-1))
\longrightarrow H^1(L(-n))\longrightarrow H^1(L_Y(-n)).
\tag{2.9.1}
\]
Por 1.7,
\(\chi(L_Y(n))=\sum_{i=0}^{r-1}a_{i+1}\binom{n+i}{i}\). Por tanto, para
\(n\ge2+h_2(a_{r-2},a_{r-1},a_r)\), la hipótesis de inducción da
\(h^1(L_Y(-n))=0\), y como \(n\ge2\) por (ii) para \(r=2\), 2.4 da
\(h^0(L_Y(-n))=0\)\,; de donde
\(h^1(L(-n-1))=h^1(L(-n))\). Pero, para \(n\gg0\),
\(h^1(L(\sisi{n}{-n}))=0\) (véase FGA 149 \S 6)\,; de donde resulta (i). % Canon SGA-SRC-001430

Como hipótesis de inducción, se supone que los polinomios \(g_i\) y \(h_i\)
están definidos para \(i\le r-1\). Entonces, para todo \(n\),
\(h^1(L(n))\le h^1(L(n-1))+h_{r-1}(a_{r-2},a_{r-1},a_r)\), gracias a
2.9.1. Por tanto, por (i), ya demostrado, se concluye que, cualquiera que sea
\(m\ge -2\),
\[
h^1(L(n))\le
\bigl[2+h_2(a_{r-2},a_{r-1},a_r)+m\bigr]
h_{r-1}(a_{r-2},a_{r-1},a_r)
\tag{2.9.2}
\]
para todo \(n\le m\).

Por 2.5, \(L_Y\) es un \((0,\ldots,0,a_r)\)-haz. Sean
\(c_i=g_{r-i}(a_{r-2},a_{r-1},a_r)\) para \(i=1,\ldots,r\) y
\(m=P_{r-2}(c_1,\ldots,c_{r-1})\). Por 1.11, \(m\ge0\) y \(L_Y\) es
\(m\)-regular. Por tanto, por 1.4 (ii),
\(h^1(L(m-2))\ge h^1(L(n))\) para todo \(n\ge m-1\). Así, si se definen
polinomios mediante \(M=P_{r-2}(g_{r-1},\ldots,g_1)\) y
\(h_r=(h_2+M)h_{r-1}\), se obtiene (ii) en virtud de 2.9.2.

Finalmente, como \(L_Y\) es \(m\)-regular, \(h^q(L(m))=0\) para \(q\ge2\)
por 1.4 (i)\,; de donde
\(\sum_{i=0}^{r}a_i\binom{m+i}{i}=h^0(L(m))-h^1(L(m))\). Pero
\(h^0(L(m))\le a_r\binom{m+r}{r}\) por 1.8 (i). Por tanto, si se define
\(g_r=\sum_{i=1}^{r}g_{r-i}\binom{M+i}{i}+a_r\binom{M+r}{r}+h_r\), se
obtiene (iii).

\textbf{Observación 2.10.} Para \(m\) fijo, consideremos una igualdad de
polinomios en \(n\):
\[
\sum_{i=0}^{r}a_{m,i}\binom{n+i}{i}
=
\sum_{j=0}^{r}a_j\binom{m+n+j}{j}.
\]
Entonces se obtiene fácilmente la fórmula
\[
a_{m,i}=\sum_{j=0}^{r-i}a_{j+i}\binom{m-1+\sisi{i}{j}}{j} % Canon SGA-SRC-001431
\]
razonando por inducción sobre \(r\) para hacer que \(a_{m,i}\) se convierta
en \(a_{m,0}\), y tomando \(n=-1\).

Del mismo modo, se obtiene que existe un polinomio con coeficientes racionales
\(P_i^{(m,r)}(X_r,\ldots,X_i)\), creciente en \(X_i\), tal que, en la
igualdad siguiente:
\[
\sum_{i=0}^{r}a_i^{(m)}\binom{n+i}{i}
=
\sum_{j=0}^{r}a_j\binom{mn+j}{j},
\]
se tiene: \(a_i^{(m)}=P_i^{(m,r)}(a_r,\ldots,a_i)\).

\textbf{Lema 2.11.} \emph{Sean \(S\) un esquema noetheriano y \(X\) un
\(S\)-esquema proyectivo, provisto de un haz muy amplio
\(H=\mathcal O_X(1)\). Supongamos que las fibras de \(X/S\) son
geométricamente íntegras, de la misma dimensión \(r\). Sea \(\mathcal L\)
una familia de clases de haces invertibles \(L_K\) sobre las fibras de
\(X/S\). Entonces, para que \(\mathcal L\) sea acotada, (es necesario y)
basta con que, en los polinomios de Hilbert
\(\chi(L_K(n))=\sum_{i=0}^{r}a_i\binom{n+i}{i}\) de los
\(L_K\in\mathcal L\), todos los coeficientes \(a_i\) permanezcan acotados.
Además, si las fibras de \(X/S\) son también geométricamente normales, basta
con que los tres coeficientes \(a_r,a_{r-1},a_{r-2}\) permanezcan acotados.}

En efecto, sea \(\delta=\sisi{\operatorname{supp}}{\sup}\{d(L_K)\}\). Si \(\delta>0\), % Canon SGA-SRC-001432
se puede sustituir \(\mathcal L\) por la familia de las clases de los
\(L_K(-\delta)\), gracias a 2.3 y 2.10\,; se puede suponer así que
\(\delta\le0\). Entonces \(\mathcal L\) es una \((b)\)-familia según 2.5,
por lo que es acotada si todos los coeficientes \(a_i\) permanecen acotados,
en virtud de 1.13; y, cuando las fibras de \(X/S\) son geométricamente
normales, basta para ello con que los tres coeficientes
\(a_r,a_{r-1},a_{r-2}\) permanezcan acotados, en virtud de 2.9.

\subsection*{3. Teoremas generales de finitud}
\textbf{Lema 3.1.} \emph{Sean \(S\) un esquema y \(X\) un \(S\)-esquema
cuasicompacto y cuasiseparado; sea \(\underline{\Pic}_{X/S}\) el funtor de
Picard (localizado para la topología f.p.p.f.)}\footnote{Cf. SGA 3 IV 6.3 y
FGA n.º 232 para la definición de \(\underline{\Pic}_{X/S}\), donde
conviene aquí, como se ha indicado, sustituir la topología fpqc por la
topología fppf.} \emph{Entonces, para todo sistema proyectivo filtrante
\((Z_\alpha)\) de \(S\)-esquemas cuasicompactos y cuasiseparados, con
morfismos de transición afines, la aplicación canónica}
\[
\varinjlim_\alpha \underline{\Pic}_{X/S}(Z_\alpha)\longrightarrow
\underline{\Pic}_{X/S}\bigl(\varprojlim_\alpha Z_\alpha\bigr)
\tag{3.1.1}
\]
\emph{es biyectiva.}

\emph{Por consiguiente} (EGA IV 8.14.2), \emph{si el funtor de Picard está
representado por un \(S\)-esquema \(\Pic_{X/S}\), entonces \(\Pic_{X/S}\) es
localmente de presentación finita.}

En efecto, un elemento \(\lambda_\alpha\) de
\(\underline{\Pic}_{X/S}(Z_\alpha)\) está representado por un haz invertible
\(L'_\alpha\) sobre un \(X_{Z'_\alpha}\), donde
\(Z'_\alpha\to Z_\alpha\) es un morfismo f.p.p.f. Pongamos
\(Z=\varprojlim Z_\alpha\). La imagen \(\lambda\) de \(\lambda_\alpha\) en
\(\underline{\Pic}_{X/S}(Z)\) está representada por el haz invertible \(L'\)
sobre \(X_{Z'}\), imagen inversa de \(\sisi{L_\alpha}{L'_\alpha}\), donde % Canon SGA-SRC-001433
\(Z'=Z'_\alpha\times_{Z_\alpha}Z\).

Si \(\lambda\) es el elemento neutro de \(\underline{\Pic}_{X/S}(Z)\),
existe un morfismo f.p.p.f. \(Z''\to Z'\) tal que la imagen inversa de \(L'\)
sobre \(X_{Z''}\) sea \(\simeq\mathcal O_{X_{Z''}}\). En virtud de EGA IV
8.8.2, 8.10.5 (vi), 11.2.6 (ii), 8.5.2 (i) y 8.5.2.4, existe un
\(\beta\geq\alpha\) y un morfismo f.p.p.f.
\(Z''_\beta\to Z'_\beta\), donde
\(Z'_\beta=Z'_\alpha\times_{Z_\alpha}Z_\beta\), tal que la imagen inversa de
\(L'_\alpha\) sobre \(X_{Z''_\beta}\) sea
\(\simeq\mathcal O_{X_{Z''_\beta}}\). La imagen de \(\lambda_\alpha\) ya es,
por tanto, el elemento neutro en \(\underline{\Pic}_{X/S}(Z_\beta)\). Por
consiguiente, 3.1.1 es inyectiva.

La sobreyectividad de 3.1.1 se demuestra de manera análoga.

\textbf{Proposición 3.2.} \emph{Sean \(S\) un esquema noetheriano y \(X\) un
\(S\)-esquema propio. Supongamos que el funtor de Picard está representado
por un \(S\)-esquema \(\Pic_{X/S}\). Entonces}:
\begin{enumerate}[label=(\roman*)]
\item \emph{\(\Pic_{X/S}\) es localmente de tipo finito sobre \(S\).}
\item \emph{Los subconjuntos \(\Lambda\) de \(\Pic_{X/S}\) corresponden
biyectivamente a las familias \(\mathscr L\) de clases de haces invertibles
sobre las fibras de \(X/S\).}
\item \emph{Una familia \(\Lambda\) de \(\Pic_{X/S}\) es cuasicompacta si y
solo si la familia correspondiente \(\mathscr L\) es acotada.}
\end{enumerate}

En efecto, (i) resulta de (3.1); (ii), del hecho siguiente: si \(s\) es un
punto de \(S\) y \(K\) es una extensión algebraicamente cerrada de \(k(s)\),
entonces los \(K\)-puntos de \(\Pic_{X/S}\) corresponden biyectivamente a
las clases de isomorfismo de haces invertibles sobre \(X_K\). En cuanto a
(iii), supongamos que el subconjunto \(\Lambda\) de \(\Pic_{X/S}\) es
cuasicompacto. Entonces \(\Lambda\) está contenido en un abierto \(U\) de
tipo finito, y la inclusión \(U\to\Pic_{X/S}\) corresponde a un haz
invertible \(L\) sobre un \(X_T\), donde \(T\to U\) es un morfismo
sobreyectivo, (plano) de tipo finito. Finalmente, \(T\to S\) es de tipo
finito, por lo que \(L\) acota \(\mathscr L\). Recíprocamente, supongamos que
\(\mathscr L\) está acotada por un haz coherente \(L\) sobre un \(X_T\), con
\(T\) de tipo finito sobre \(S\). Se puede suponer que \(L\) es invertible
(1.13 (i)), de modo que \(L\) defina un morfismo
\(f:T\to\Pic_{X/S}\). La imagen de \(f\) contiene a \(\Lambda\) y es
cuasicompacta, luego es noetheriana por (i). Por consiguiente, \(\Lambda\) es
cuasicompacto y se obtiene (iii).

Gracias a 3.2, la definición habitual de los morfismos de tipo finito entre
dos esquemas de Picard se generaliza en la siguiente:

\textbf{Definición 3.3.} \emph{Sean \(X\) e \(Y\) propios sobre \(S\)
noetheriano. Un morfismo entre los funtores de Picard}
\[
\begin{tikzcd}
\underline{\Pic}_{Y/S}\arrow[r]&\underline{\Pic}_{X/S}
\end{tikzcd}
\]
\emph{se dice de tipo finito si toda familia acotada de clases de haces
invertibles sobre las fibras de \(X/S\) tiene como imagen inversa una familia
acotada de clases de haces sobre las fibras de \(Y/S\).}

\textbf{Lema 3.4.} \emph{Sean \(X,Y,X',Y'\) propios sobre \(S\) noetheriano
y \(T\to S\) sobreyectivo de tipo finito. Entonces}:
\begin{enumerate}[label=(\roman*)]
\item \emph{Para que un morfismo
\(\Phi:\underline{\Pic}_{Y/S}\to\underline{\Pic}_{X/S}\) sea de tipo finito,
es necesario y suficiente que
\(\Phi_T:\underline{\Pic}_{Y_T/T}\to\underline{\Pic}_{X_T/T}\) sea de tipo
finito.}
\item \emph{En un diagrama conmutativo}
\[
\begin{tikzcd}
\Phi':\underline{\Pic}_{Y'/S}\arrow[r]&
\underline{\Pic}_{X'/S}\\
